% GENERATED FILE. Edit canonical lessons, then run npm run generate:books.
% canonical-source-hash: fnv1a64:b31e355d73c99611
% canonical-lessons: MR-C490-pathya, MR-C490-nirogi, MR-C490-pratikarsakti, MR-C490-sekne, MR-C490-colne

\chapter{Staying Healthy}
\label{ch:mr-staying-healthy}
\hlchaptermodality{490}

\begin{chapteropening}
\textbf{By the end of this chapter:} \emph{I can talk about a diet regimen, good health, immunity and home remedies.}

Complete the last lesson of chapter 490: I can talk about a diet regimen, good health, immunity and home remedies.
\end{chapteropening}

\section[\texorpdfstring{\mr{पथ्य} (pathya)}{पथ्य (pathya)}]{\mr{पथ्य} (pathya) — a diet regimen}

\label{lesson:MR-C490-pathya}



\begin{quote}
Before the new one: say the Marathi for a stitch, then the Marathi for to prick.
\end{quote}

\subsection*{You'll want to know: \mr{पथ्य}}
\textbf{\mr{पथ्य}} — \emph{pathya} — \textquotedblleft{}a diet regimen\textquotedblright{}.

\begin{quote}
\mr{आजारपणात} \mr{आजीने} \mr{कडक} \mr{पथ्य} \mr{पाळलं}.
\end{quote}

\begin{quote}
\emph{ājārpaṇāt ājīne kaḍak pathya pāḷlaṁ.}
\end{quote}

\begin{quote}
During her illness Grandmother kept to a strict diet.
\end{quote}

Neuter.

\mr{पथ्य} \mr{पाळणे}, to keep to a prescribed diet; its opposite is \mr{कुपथ्य}.

\begin{cousinweb}
\textbf{\mr{पथ्य}} (\emph{pathya}) is a tatsama, taken from Sanskrit \emph{pathya}, \textquotedblleft{}suitable\textquotedblright{} (CDIAL 7746): what suits the patient.
\end{cousinweb}

\begin{grammarlens}[title={What you've built}]
One more word for describing what you lived and felt.
\end{grammarlens}

\subsection*{Your turn}
\begin{itemize}
\raggedright
  \item \emph{Say it:} \emph{pathya}
  \item \emph{Say it:} \emph{pathya}, once more
  \item \emph{Say it:} \emph{ājārpaṇāt ājīne kaḍak pathya pāḷlaṁ.} — During her illness Grandmother kept to a strict diet.
  \item \emph{Recall:} say the Marathi for a stitch, then the Marathi for to prick, then say \emph{pathya} again
\end{itemize}

\begin{tcolorbox}[breakable,colback=teal!4,colframe=teal!35!black,title={Before you move on}]
What does \mr{पथ्य} mean? (\textquotedblleft{}A diet regimen\textquotedblright{}.) Say it once more.
\end{tcolorbox}
\section[\texorpdfstring{\mr{निरोगी} (nirogī)}{निरोगी (nirogī)}]{\mr{निरोगी} (nirogī) — healthy}

\label{lesson:MR-C490-nirogi}



\begin{quote}
Before the new one: say the Marathi for to prick, then the Marathi for a diet regimen.
\end{quote}

\subsection*{You'll want to know: \mr{निरोगी}}
\textbf{\mr{निरोगी}} — \emph{nirogī} — \textquotedblleft{}healthy\textquotedblright{}.

\begin{quote}
\mr{रोज} \mr{चालल्याने} \mr{आपण} \mr{निरोगी} \mr{राहतो}.
\end{quote}

\begin{quote}
\emph{roj cāllyāne āpaṇ nirogī rāhto.}
\end{quote}

\begin{quote}
Walking every day keeps us healthy.
\end{quote}

\mr{नि}-, without, + \mr{रोग}, disease, + -\mr{ई}. Unchanging.

\begin{cousinweb}
\textbf{\mr{निरोगी}} (\emph{nirogī}) is a Marathi form built on Sanskrit \emph{nīroga}, \textquotedblleft{}free from sickness, healthy\textquotedblright{} (Monier-Williams), with the adjective ending \emph{-ī}.
\end{cousinweb}

\begin{grammarlens}[title={What you've built}]
One more word for describing what you lived and felt.
\end{grammarlens}

\subsection*{Your turn}
\begin{itemize}
\raggedright
  \item \emph{Say it:} \emph{nirogī}
  \item \emph{Say it:} \emph{nirogī}, once more
  \item \emph{Say it:} \emph{roj cāllyāne āpaṇ nirogī rāhto.} — Walking every day keeps us healthy.
  \item \emph{Recall:} say the Marathi for to prick, then the Marathi for a diet regimen, then say \emph{nirogī} again
\end{itemize}

\begin{tcolorbox}[breakable,colback=teal!4,colframe=teal!35!black,title={Before you move on}]
What does \mr{निरोगी} mean? (\textquotedblleft{}Healthy\textquotedblright{}.) Say it once more.
\end{tcolorbox}
\section[\texorpdfstring{\mr{प्रतिकारशक्ती} (pratikārśaktī)}{प्रतिकारशक्ती (pratikārśaktī)}]{\mr{प्रतिकारशक्ती} (pratikārśaktī) — immunity; resistance}

\label{lesson:MR-C490-pratikarsakti}



\begin{quote}
Before the new one: say the Marathi for a diet regimen, then the Marathi for healthy.
\end{quote}

\subsection*{You'll want to know: \mr{प्रतिकारशक्ती}}
\textbf{\mr{प्रतिकारशक्ती}} — \emph{pratikārśaktī} — \textquotedblleft{}immunity; resistance\textquotedblright{}.

\begin{quote}
\mr{चांगल्या} \mr{आहारामुळे} \mr{प्रतिकारशक्ती} \mr{वाढते}.
\end{quote}

\begin{quote}
\emph{cāṅglyā āhārāmuḷe pratikārśaktī vāḍhte.}
\end{quote}

\begin{quote}
A good diet builds up immunity.
\end{quote}

Feminine.

\mr{प्रतिकार}, resistance, + \mr{शक्ती}, power.

\begin{cousinweb}
\textbf{\mr{प्रतिकारशक्ती}} (\emph{pratikārśaktī}) is a modern learned compound built from two Sanskrit words, \emph{pratikāra}, \textquotedblleft{}requital, counteraction\textquotedblright{}, and \emph{śakti}, \textquotedblleft{}power\textquotedblright{} (Monier-Williams).
\end{cousinweb}

\begin{grammarlens}[title={What you've built}]
One more word for describing what you lived and felt.
\end{grammarlens}

\subsection*{Your turn}
\begin{itemize}
\raggedright
  \item \emph{Say it:} \emph{pratikārśaktī}
  \item \emph{Say it:} \emph{pratikārśaktī}, once more
  \item \emph{Say it:} \emph{cāṅglyā āhārāmuḷe pratikārśaktī vāḍhte.} — A good diet builds up immunity.
  \item \emph{Recall:} say the Marathi for a diet regimen, then the Marathi for healthy, then say \emph{pratikārśaktī} again
\end{itemize}

\begin{tcolorbox}[breakable,colback=teal!4,colframe=teal!35!black,title={Before you move on}]
What does \mr{प्रतिकारशक्ती} mean? (\textquotedblleft{}Immunity; resistance\textquotedblright{}.) Say it once more.
\end{tcolorbox}
\section[\texorpdfstring{\mr{शेकणे} (śekṇe)}{शेकणे (śekṇe)}]{\mr{शेकणे} (śekṇe) — to apply heat; to warm}

\label{lesson:MR-C490-sekne}



\begin{quote}
Before the new one: say the Marathi for healthy, then the Marathi for immunity.
\end{quote}

\subsection*{You'll want to know: \mr{शेकणे}}
\textbf{\mr{शेकणे}} — \emph{śekṇe} — \textquotedblleft{}to apply heat; to warm\textquotedblright{}.

\begin{quote}
\mr{आजीने} \mr{गरम} \mr{पाण्याच्या} \mr{पिशवीने} \mr{पाठ} \mr{शेकली}.
\end{quote}

\begin{quote}
\emph{ājīne garam pāṇyācyā piśvīne pāṭh śeklī.}
\end{quote}

\begin{quote}
Grandmother warmed her back with a hot-water bag.
\end{quote}

Past, with -\mr{ने} (he …): \mr{त्याने} \mr{पाठ} \mr{शेकली}, \emph{tyāne pāṭh śeklī}.

Gentle heat on a sore spot; \mr{शेक}, a warm compress.

\begin{cousinweb}
\textbf{\mr{शेकणे}} (\emph{śekṇe}) is inherited: CDIAL lists Marathi \emph{śekṇẽ}, \textquotedblleft{}to warm oneself before a fire, foment\textquotedblright{}, under the reconstructed \emph{*sēkk} (CDIAL 13581).
\end{cousinweb}

\begin{grammarlens}[title={What you've built}]
One more word for describing what you lived and felt.
\end{grammarlens}

\subsection*{Your turn}
\begin{itemize}
\raggedright
  \item \emph{Say it:} \emph{śekṇe}
  \item \emph{Say it:} \emph{śekṇe}, once more
  \item \emph{Say it:} \emph{ājīne garam pāṇyācyā piśvīne pāṭh śeklī.} — Grandmother warmed her back with a hot-water bag.
  \item \emph{Recall:} say the Marathi for healthy, then the Marathi for immunity, then say \emph{śekṇe} again
\end{itemize}

\begin{tcolorbox}[breakable,colback=teal!4,colframe=teal!35!black,title={Before you move on}]
What does \mr{शेकणे} mean? (\textquotedblleft{}To apply heat; to warm\textquotedblright{}.) Say it once more.
\end{tcolorbox}
\section[\texorpdfstring{\mr{चोळणे} (coḷṇe)}{चोळणे (coḷṇe)}]{\mr{चोळणे} (coḷṇe) — to rub}

\label{lesson:MR-C490-colne}



\begin{quote}
Before the new one: say the Marathi for immunity, then the Marathi for to apply heat.
\end{quote}

\subsection*{You'll want to know: \mr{चोळणे}}
\textbf{\mr{चोळणे}} — \emph{coḷṇe} — \textquotedblleft{}to rub\textquotedblright{}.

\begin{quote}
\mr{थंडीत} \mr{त्याने} \mr{हात} \mr{चोळले}.
\end{quote}

\begin{quote}
\emph{thaṇḍīt tyāne hāt coḷle.}
\end{quote}

\begin{quote}
He rubbed his hands in the cold.
\end{quote}

Past, with -\mr{ने} (he …): \mr{त्याने} \mr{हात} \mr{चोळले}, \emph{tyāne hāt coḷle}.

Of oil or a sore spot: \mr{तेल} \mr{चोळणे}, to rub in oil.

\begin{cousinweb}
\textbf{\mr{चोळणे}} (\emph{coḷṇe}) is a Marathi verb whose origin is not settled.
\end{cousinweb}

\begin{grammarlens}[title={What you've built}]
One more word for describing what you lived and felt.
\end{grammarlens}

\subsection*{Your turn}
\begin{itemize}
\raggedright
  \item \emph{Say it:} \emph{coḷṇe}
  \item \emph{Say it:} \emph{coḷṇe}, once more
  \item \emph{Say it:} \emph{thaṇḍīt tyāne hāt coḷle.} — He rubbed his hands in the cold.
  \item \emph{Recall:} say the Marathi for immunity, then the Marathi for to apply heat, then say \emph{coḷṇe} again
\end{itemize}

\begin{tcolorbox}[breakable,colback=teal!4,colframe=teal!35!black,title={Before you move on}]
What does \mr{चोळणे} mean? (\textquotedblleft{}To rub\textquotedblright{}.) Say it once more.
\end{tcolorbox}
